\documentclass{article}
\usepackage{amsmath,amssymb,amsthm}
\IfFileExists{hyperref.sty}{%
  \IfFileExists{color.sty}{\usepackage{color}\definecolor{paperlinkblue}{rgb}{0.10,0.25,0.42}%
    \usepackage[colorlinks=true,linkcolor=paperlinkblue]{hyperref}}{\usepackage[hidelinks]{hyperref}}%
}{}
\ifdefined\hyperlink
  \newcommand{\paperanchor}[1]{\hypertarget{#1}{}}%
  \newcommand{\paperlink}[2]{\hyperlink{#1}{\begin{NoHyper}#2\end{NoHyper}}}%
\else
  \ifdefined\pdfextension
    \ifnum\outputmode>0
      \def\paperdest{\pdfextension dest}\def\paperstart{\pdfextension startlink}%
      \def\paperend{\pdfextension endlink}\def\papercolor{\pdfextension colorstack}%
    \fi
  \else\ifdefined\pdfstartlink
    \ifnum\pdfoutput>0
      \let\paperdest\pdfdest\let\paperstart\pdfstartlink\let\paperend\pdfendlink\let\papercolor\pdfcolorstack
    \fi
  \fi\fi
  \ifdefined\paperdest
    \newcommand{\paperanchor}[1]{\paperdest name{#1} xyz}%
    \newcommand{\paperlink}[2]{\leavevmode\paperstart attr{/Border [0 0 0]} goto name{#1}%
      \papercolor 0 push {0.10 0.25 0.42 rg 0.10 0.25 0.42 RG}#2\papercolor 0 pop\paperend}%
  \else\ifdefined\XeTeXversion
    \newcommand{\paperanchor}[1]{\special{pdf:dest (#1) [@thispage /XYZ null @ypos null]}}%
    \newcommand{\paperlink}[2]{\leavevmode\special{pdf:bann << /Subtype /Link /Border [0 0 0] /A << /S /GoTo /D (#1) >> >>}%
      \special{color push rgb 0.10 0.25 0.42}#2\special{color pop}\special{pdf:eann}}%
  \else
    \PackageError{lean-to-latex}{No PDF hyperlink backend available}{Use XeLaTeX, pdfLaTeX, LuaLaTeX, or install hyperref.}%
  \fi\fi
\fi
\newcommand{\paperlineanchor}[1]{\raisebox{10pt}[0pt][0pt]{\paperanchor{#1}}}
\DeclareRobustCommand{\paperresultname}[1]{\paperlineanchor{theorem:\arabic{theorem}}#1}
\newtheoremstyle{paperresult}{8pt}{5pt}{\normalfont\raggedright}{}{\bfseries}{.}{\newline}{}
\theoremstyle{paperresult}
\newtheorem{theorem}{\paperresultname{Theorem}}
\newcounter{paperpremise}
\newtheorem{lemma}[theorem]{\paperresultname{Lemma}}
\newtheorem{corollary}[theorem]{\paperresultname{Corollary}}
\newcounter{definingequation}
\newenvironment{proofstep}{\par\smallskip\ifdim\leftskip<3em\relax\advance\leftskip by 1em\relax\fi\noindent\ignorespaces}{\par\smallskip}
\setlength{\textwidth}{6.5in}
\setlength{\oddsidemargin}{0in}
\AddToHook{env/proof/begin}{\raggedright}
\setlength{\emergencystretch}{\linewidth}
\title{Double counting in finite relations}
\author{}
\date{}
\begin{document}
\maketitle
\section*{Definitions}
\paragraph{Bipartite below.} \par\noindent\refstepcounter{definingequation}\label{defeq:1}\paperlineanchor{defeq:1}\textup{(\paperlink{defeq:1}{D1})} Let \(\alpha\) and \(\beta\) be types. Let \(r\) be a relation from \(\alpha\) to \(\beta\). Let \(s\) be a finite subset of \(\alpha\). Let \(b\) be an element of \(\beta\). Fix a structure \(\mathit{inst}\) of type \(\Pi_{a\, :\, \alpha}\, \operatorname{Decidable}\left(r\left(a, b\right)\right)\). Then \(N^{-}_{r, s}(b) = \{a\in s \mid r\left(a, b\right)\}\).

\paragraph{Bipartite above.} \par\noindent\refstepcounter{definingequation}\label{defeq:2}\paperlineanchor{defeq:2}\textup{(\paperlink{defeq:2}{D2})} Let \(\alpha\) and \(\beta\) be types. Let \(r\) be a relation from \(\alpha\) to \(\beta\). Let \(t\) be a finite subset of \(\beta\). Let \(a\) be an element of \(\alpha\). Assume the predicate \(r\left(a\right)\) is decidable. Then \(N^{+}_{r, t}(a) = \{b\in t \mid r\left(a, b\right)\}\).

\subsection*{Prerequisites}
\begin{description}
\item[\paperlineanchor{premise:45712e6765}\paperlink{premise:45712e6765}{P1.}] \refstepcounter{paperpremise}\label{premise:45712e6765}\textbf{Equality implies the reverse weak inequality.} Let \(\alpha\) be a type. Equip \(\alpha\) with a preorder. Let \(a\) and \(b\) be elements of \(\alpha\). Assume \(a = b\). Then \(b \le a\).
\item[\paperlineanchor{premise:45712e6c65}\paperlink{premise:45712e6c65}{P2.}] \refstepcounter{paperpremise}\label{premise:45712e6c65}\textbf{Alias.} Let \(\alpha\) be a type. Equip \(\alpha\) with a preorder. Let \(a\) and \(b\) be elements of \(\alpha\). Assume \(a = b\). Then \(a \le b\).
\item[\paperlineanchor{premise:46696e7365742e636172645f65715f73756d5f6f6e6573}\paperlink{premise:46696e7365742e636172645f65715f73756d5f6f6e6573}{P3.}] \refstepcounter{paperpremise}\label{premise:46696e7365742e636172645f65715f73756d5f6f6e6573}\textbf{Cardinality as a sum of ones.} Let \(\iota\) be a type. Let \(s\) be a finite subset of \(\iota\). Then \(\left|s\right| = \sum_{x\in s}\left(1\right)\).
\item[\paperlineanchor{premise:46696e7365742e636172645f6c655f6f6e65}\paperlink{premise:46696e7365742e636172645f6c655f6f6e65}{P4.}] \refstepcounter{paperpremise}\label{premise:46696e7365742e636172645f6c655f6f6e65}\textbf{Cardinality at most one.} Let \(\alpha\) be a type. Let \(s\) be a finite subset of \(\alpha\). Then \(|s| \le 1 \iff (\forall a\, :\, \alpha,\; a \in s \implies (\forall b\, :\, \alpha,\; b \in s \implies a = b))\).
\item[\paperlineanchor{premise:46696e7365742e636172645f6e736d756c5f6c655f73756d}\paperlink{premise:46696e7365742e636172645f6e736d756c5f6c655f73756d}{P5.}] \refstepcounter{paperpremise}\label{premise:46696e7365742e636172645f6e736d756c5f6c655f73756d}\textbf{Cardinality repeated addition at most sum.} Let \(\iota\) and \(N\) be types. Equip \(N\) with a commutative additive monoid structure and a partial order. Assume \(\text{translation preserves order on }N\). Let \(s\) be a finite subset of \(\iota\). Let \(f\) be a function of type \(\iota \to N\). Let \(n\) be an element of \(N\). Assume \(\forall x\, :\, \iota,\; x \in s \implies n \le f\left(x\right)\). Then \(\sum_{k=1}^{\left|s\right|}\left(n\right) \le \sum_{i\in s}\left(f(i)\right)\).
\item[\paperlineanchor{premise:46696e7365742e636172645f706f73}\paperlink{premise:46696e7365742e636172645f706f73}{P6.}] \refstepcounter{paperpremise}\label{premise:46696e7365742e636172645f706f73}\textbf{Cardinality positivity.} Let \(\alpha\) be a type. Let \(s\) be a finite subset of \(\alpha\). Then \(0 < \left|s\right| \iff s\ne\varnothing\).
\item[\paperlineanchor{premise:46696e7365742e636f655f66696c746572}\paperlink{premise:46696e7365742e636f655f66696c746572}{P7.}] \refstepcounter{paperpremise}\label{premise:46696e7365742e636f655f66696c746572}\textbf{Coercion filter.} Let \(\alpha\) be a type. Let \(p\) be a function of type \(\alpha \to \mathsf{Prop}\). Assume the predicate \(p\) is decidable. Let \(s\) be a finite subset of \(\alpha\). Then \(\operatorname{underlying\,set}(\{x\in s \mid p(x)\}) = \{x\in \alpha \mid x \in s \land p\left(x\right)\}\).
\item[\paperlineanchor{premise:46696e7365742e636f655f6e6f6e656d707479}\paperlink{premise:46696e7365742e636f655f6e6f6e656d707479}{P8.}] \refstepcounter{paperpremise}\label{premise:46696e7365742e636f655f6e6f6e656d707479}\textbf{Coercion nonempty.} Let \(\alpha\) be a type. Let \(s\) be a finite subset of \(\alpha\). Then \(\operatorname{underlying\,set}(s)\ne\varnothing \iff s\ne\varnothing\).
\item[\paperlineanchor{premise:46696e7365742e6d656d5f66696c746572}\paperlink{premise:46696e7365742e6d656d5f66696c746572}{P9.}] \refstepcounter{paperpremise}\label{premise:46696e7365742e6d656d5f66696c746572}\textbf{Membership filter.} Let \(\alpha\) be a type. Let \(p\) be a function of type \(\alpha \to \mathsf{Prop}\). Assume the predicate \(p\) is decidable. Let \(s\) be a finite subset of \(\alpha\). Let \(a\) be an element of \(\alpha\). Then \(a \in \{x\in s \mid p(x)\} \iff a \in s \land p\left(a\right)\).
\item[\paperlineanchor{premise:46696e7365742e6d656d5f756e6976}\paperlink{premise:46696e7365742e6d656d5f756e6976}{P10.}] \refstepcounter{paperpremise}\label{premise:46696e7365742e6d656d5f756e6976}\textbf{Membership in the full set.} Let \(\alpha\) be a type. Fix a finite enumeration of \(\alpha\). Let \(x\) be an element of \(\alpha\). Then \(x \in \alpha\).
\item[\paperlineanchor{premise:46696e7365742e70726f645f636f6d6d}\paperlink{premise:46696e7365742e70726f645f636f6d6d}{P11.}] \refstepcounter{paperpremise}\label{premise:46696e7365742e70726f645f636f6d6d}\textbf{Interchanging two finite products.} Let \(\alpha\), \(\beta\), and \(\gamma\) be types. Fix a structure \(\mathit{inst}\) of type \(\operatorname{Comm\,Monoid}\left(\beta\right)\). Let \(s\) be a finite subset of \(\gamma\). Let \(t\) be a finite subset of \(\alpha\). Let \(f\) be a function of type \(\gamma \to \alpha \to \beta\). Then \(\prod_{x\in s}(\prod_{y\in t}(f(x,\allowbreak{} y))) = \prod_{y\in t}(\prod_{x\in s}(f(x,\allowbreak{} y)))\).
\item[\paperlineanchor{premise:46696e7365742e70726f645f636f6e6772}\paperlink{premise:46696e7365742e70726f645f636f6e6772}{P12.}] \refstepcounter{paperpremise}\label{premise:46696e7365742e70726f645f636f6e6772}\textbf{Equality of products with equal factors.} Let \(\iota\) and \(M\) be types. Let \(s_{1}\) and \(s_{2}\) be finite subsets of \(\iota\). Fix a structure \(\mathit{inst}\) of type \(\operatorname{Comm\,Monoid}\left(M\right)\). Let \(f\) and \(g\) be functions of type \(\iota \to M\). Assume \(s_{1} = s_{2}\). Assume \(\forall x\, :\, \iota,\; x \in s_{2} \implies f\left(x\right) = g\left(x\right)\). Then \(\prod_{i\in s_{1}}\left(f(i)\right) = \prod_{i\in s_{2}}\left(g(i)\right)\).
\item[\paperlineanchor{premise:46696e7365742e70726f645f66696c746572}\paperlink{premise:46696e7365742e70726f645f66696c746572}{P13.}] \refstepcounter{paperpremise}\label{premise:46696e7365742e70726f645f66696c746572}\textbf{Product filter.} Let \(\iota\) and \(M\) be types. Let \(s\) be a finite subset of \(\iota\). Fix a structure \(\mathit{inst}\) of type \(\operatorname{Comm\,Monoid}\left(M\right)\). Let \(p\) be a function of type \(\iota \to \mathsf{Prop}\). Assume the predicate \(p\) is decidable. Let \(f\) be a function of type \(\iota \to M\). Then \(\prod_{a\in \{a\in s \mid p(a)\}}(f(a)) = \prod_{a\in s}(\begin{cases}f(a) &\text{if } p(a)\\ 1 &\text{otherwise}\end{cases})\).
\item[\paperlineanchor{premise:46696e7365742e73756d5f636f6d6d}\paperlink{premise:46696e7365742e73756d5f636f6d6d}{P14.}] \refstepcounter{paperpremise}\label{premise:46696e7365742e73756d5f636f6d6d}\textbf{Interchanging two finite sums.} Let \(\alpha\), \(\beta\), and \(\gamma\) be types. Equip \(\beta\) with a commutative additive monoid structure. Let \(s\) be a finite subset of \(\gamma\). Let \(t\) be a finite subset of \(\alpha\). Let \(f\) be a function of type \(\gamma \to \alpha \to \beta\). Then \(\sum_{x\in s}(\sum_{y\in t}(f(x,\allowbreak{} y))) = \sum_{y\in t}(\sum_{x\in s}(f(x,\allowbreak{} y)))\).
\item[\paperlineanchor{premise:46696e7365742e73756d5f636f6e6772}\paperlink{premise:46696e7365742e73756d5f636f6e6772}{P15.}] \refstepcounter{paperpremise}\label{premise:46696e7365742e73756d5f636f6e6772}\textbf{Equality of sums with equal summands.} Let \(\iota\) and \(M\) be types. Let \(s_{1}\) and \(s_{2}\) be finite subsets of \(\iota\). Equip \(M\) with a commutative additive monoid structure. Let \(f\) and \(g\) be functions of type \(\iota \to M\). Assume \(s_{1} = s_{2}\). Assume \(\forall x\, :\, \iota,\; x \in s_{2} \implies f\left(x\right) = g\left(x\right)\). Then \(\sum_{i\in s_{1}}\left(f(i)\right) = \sum_{i\in s_{2}}\left(g(i)\right)\).
\item[\paperlineanchor{premise:46696e7365742e73756d5f636f6e7374}\paperlink{premise:46696e7365742e73756d5f636f6e7374}{P16.}] \refstepcounter{paperpremise}\label{premise:46696e7365742e73756d5f636f6e7374}\textbf{The sum of a constant.} Let \(\iota\) and \(M\) be types. Let \(s\) be a finite subset of \(\iota\). Equip \(M\) with a commutative additive monoid structure. Let \(b\) be an element of \(M\). Then \(\sum_{x\in s}\left(b\right) = \sum_{k=1}^{\left|s\right|}\left(b\right)\).
\item[\paperlineanchor{premise:46696e7365742e73756d5f66696c746572}\paperlink{premise:46696e7365742e73756d5f66696c746572}{P17.}] \refstepcounter{paperpremise}\label{premise:46696e7365742e73756d5f66696c746572}\textbf{Sum filter.} Let \(\iota\) and \(M\) be types. Let \(s\) be a finite subset of \(\iota\). Equip \(M\) with a commutative additive monoid structure. Let \(p\) be a function of type \(\iota \to \mathsf{Prop}\). Assume the predicate \(p\) is decidable. Let \(f\) be a function of type \(\iota \to M\). Then \(\sum_{a\in \{a\in s \mid p(a)\}}(f(a)) = \sum_{a\in s}(\begin{cases}f(a) &\text{if } p(a)\\ 0 &\text{otherwise}\end{cases})\).
\item[\paperlineanchor{premise:46696e7365742e73756d5f6c655f636172645f6e736d756c}\paperlink{premise:46696e7365742e73756d5f6c655f636172645f6e736d756c}{P18.}] \refstepcounter{paperpremise}\label{premise:46696e7365742e73756d5f6c655f636172645f6e736d756c}\textbf{Sum at most cardinality repeated addition.} Let \(\iota\) and \(N\) be types. Equip \(N\) with a commutative additive monoid structure and a partial order. Assume \(\text{translation preserves order on }N\). Let \(s\) be a finite subset of \(\iota\). Let \(f\) be a function of type \(\iota \to N\). Let \(n\) be an element of \(N\). Assume \(\forall x\, :\, \iota,\; x \in s \implies f\left(x\right) \le n\). Then \(\sum_{i\in s}\left(f(i)\right) \le \sum_{k=1}^{\left|s\right|}\left(n\right)\).
\item[\paperlineanchor{premise:46696e7365742e73756d5f6c745f73756d5f6f665f6e6f6e656d707479}\paperlink{premise:46696e7365742e73756d5f6c745f73756d5f6f665f6e6f6e656d707479}{P19.}] \refstepcounter{paperpremise}\label{premise:46696e7365742e73756d5f6c745f73756d5f6f665f6e6f6e656d707479}\textbf{Sum less than sum of nonempty.} Let \(\iota\) and \(M\) be types. Equip \(M\) with a commutative additive monoid structure and a partial order. Assume \(\text{translation preserves and reflects order on }M\). Let \(f\) and \(g\) be functions of type \(\iota \to M\). Let \(s\) be a finite subset of \(\iota\). Assume \(s\ne\varnothing\). Assume \(\forall i\, :\, \iota,\; i \in s \implies f\left(i\right) < g\left(i\right)\). Then \(\sum_{i\in s}\left(f\left(i\right)\right) < \sum_{i\in s}\left(g\left(i\right)\right)\).
\item[\paperlineanchor{premise:49734f72646572656452696e672e746f49734f7264657265644164644d6f6e6f6964}\paperlink{premise:49734f72646572656452696e672e746f49734f7264657265644164644d6f6e6f6964}{P20.}] \refstepcounter{paperpremise}\label{premise:49734f72646572656452696e672e746f49734f7264657265644164644d6f6e6f6964}\textbf{To is ordered sum monoid.} Let \(R\) be a type. Let \(\mathit{inst}\) be an element of \(\operatorname{Semiring}\left(R\right)\). Let \(\mathit{inst}_{1}\) be an element of \(\operatorname{Partial\,Order}\left(R\right)\). Let \(\mathit{self}\) be a proof of \(\operatorname{Is\,Ordered\,Ring}\left(R\right)\). Then \(\text{translation preserves order on }R\).
\item[\paperlineanchor{premise:49735374726963744f72646572656452696e672e746f436861725a65726f}\paperlink{premise:49735374726963744f72646572656452696e672e746f436861725a65726f}{P21.}] \refstepcounter{paperpremise}\label{premise:49735374726963744f72646572656452696e672e746f436861725a65726f}\textbf{To char zero.} Let \(R\) be a type. Equip \(R\) with a semiring structure. Equip \(R\) with a partial order. Assume \(R\text{ is a strictly ordered semiring}\). Then \(\operatorname{Char\,Zero}\left(R\right)\).
\item[\paperlineanchor{premise:49735374726963744f72646572656452696e672e746f49734f72646572656443616e63656c4164644d6f6e6f6964}\paperlink{premise:49735374726963744f72646572656452696e672e746f49734f72646572656443616e63656c4164644d6f6e6f6964}{P22.}] \refstepcounter{paperpremise}\label{premise:49735374726963744f72646572656452696e672e746f49734f72646572656443616e63656c4164644d6f6e6f6964}\textbf{To is ordered cancel sum monoid.} Let \(R\) be a type. Let \(\mathit{inst}\) be an element of \(\operatorname{Semiring}\left(R\right)\). Let \(\mathit{inst}_{1}\) be an element of \(\operatorname{Partial\,Order}\left(R\right)\). Assume \(R\text{ is a strictly ordered semiring}\). Then \(\text{translation preserves and reflects order on }R\).
\item[\paperlineanchor{premise:49735374726963744f72646572656452696e672e746f49734f72646572656452696e67}\paperlink{premise:49735374726963744f72646572656452696e672e746f49734f72646572656452696e67}{P23.}] \refstepcounter{paperpremise}\label{premise:49735374726963744f72646572656452696e672e746f49734f72646572656452696e67}\textbf{To is ordered ring.} Let \(R\) be a type. Equip \(R\) with a semiring structure. Equip \(R\) with a partial order. Assume \(R\text{ is a strictly ordered semiring}\). Then \(\operatorname{Is\,Ordered\,Ring}\left(R\right)\).
\item[\paperlineanchor{premise:4c452e6c652e616e746973796d6d}\paperlink{premise:4c452e6c652e616e746973796d6d}{P24.}] \refstepcounter{paperpremise}\label{premise:4c452e6c652e616e746973796d6d}\textbf{Alias.} Let \(\alpha\) be a type. Equip \(\alpha\) with a partial order. Let \(a\) and \(b\) be elements of \(\alpha\). Assume \(a \le b\). Assume \(b \le a\). Then \(a = b\).
\item[\paperlineanchor{premise:4e61742e636173745f696e6a}\paperlink{premise:4e61742e636173745f696e6a}{P25.}] \refstepcounter{paperpremise}\label{premise:4e61742e636173745f696e6a}\textbf{Cast injectivity.} Let \(R\) be a type. Fix a structure \(\mathit{inst}\) of type \(\operatorname{Add\,Monoid\,With\,One}\left(R\right)\). Let \(\mathit{inst}_{1}\) be a proof of \(\operatorname{Char\,Zero}\left(R\right)\). Let \(m\) and \(n\) be natural numbers. Then \(\operatorname{cast}_{R}\left(m\right) = \operatorname{cast}_{R}\left(n\right) \iff m = n\).
\item[\paperlineanchor{premise:4e61742e636173745f73756d}\paperlink{premise:4e61742e636173745f73756d}{P26.}] \refstepcounter{paperpremise}\label{premise:4e61742e636173745f73756d}\textbf{Cast sum.} Let \(\alpha\) and \(\beta\) be types. Fix a structure \(\mathit{inst}\) of type \(\operatorname{Add\,Comm\,Monoid\,With\,One}\left(\beta\right)\). Let \(s\) be a finite subset of \(\alpha\). Let \(f\) be a function of type \(\alpha \to \mathbb{N}\). Then \(\operatorname{cast}_{\beta}(\sum_{x\in s}(f(x))) = \sum_{x\in s}(\operatorname{cast}_{\beta}(f(x)))\).
\item[\paperlineanchor{premise:4e61742e696e737449735374726963744f72646572656452696e67}\paperlink{premise:4e61742e696e737449735374726963744f72646572656452696e67}{P27.}] \refstepcounter{paperpremise}\label{premise:4e61742e696e737449735374726963744f72646572656452696e67}\textbf{Strict order on the natural numbers.} \(\mathbb{N}\text{ is a strictly ordered semiring}\).
\item[\paperlineanchor{premise:65715f74727565}\paperlink{premise:65715f74727565}{P28.}] \refstepcounter{paperpremise}\label{premise:65715f74727565}\textbf{Equals true.} Let \(p\) be a proposition. Assume \(p\). Then \(p = \mathrm{True}\).
\item[\paperlineanchor{premise:666f72616c6c5f636f6e6772}\paperlink{premise:666f72616c6c5f636f6e6772}{P29.}] \refstepcounter{paperpremise}\label{premise:666f72616c6c5f636f6e6772}\textbf{Forall congruence.} Let \(\alpha\) be a sort. Let \(p\) and \(q\) be functions of type \(\alpha \to \mathsf{Prop}\). Assume \(\forall a\, :\, \alpha,\; p\left(a\right) = q\left(a\right)\). Then \((\forall a\, :\, \alpha,\; p(a)) = (\forall a\, :\, \alpha,\; q(a))\).
\item[\paperlineanchor{premise:666f72616c6c5f636f6e7374}\paperlink{premise:666f72616c6c5f636f6e7374}{P30.}] \refstepcounter{paperpremise}\label{premise:666f72616c6c5f636f6e7374}\textbf{Forall constant.} Let \(b\) be a proposition. Let \(\alpha\) be a sort. Assume \(\alpha\text{ is nonempty}\). Then \(\left(\forall a\, :\, \alpha,\; b\right) \iff b\).
\item[\paperlineanchor{premise:696d706c6965735f636f6e6772}\paperlink{premise:696d706c6965735f636f6e6772}{P31.}] \refstepcounter{paperpremise}\label{premise:696d706c6965735f636f6e6772}\textbf{Implies congruence.} Let \(p_{1}\) and \(p_{2}\) be sorts. Let \(q_{1}\) and \(q_{2}\) be sorts. Assume \(p_{1} = p_{2}\). Assume \(q_{1} = q_{2}\). Then \(\left(p_{1} \to q_{1}\right) = \left(p_{2} \to q_{2}\right)\).
\item[\paperlineanchor{premise:696e73744e6f6e656d7074794f66496e68616269746564}\paperlink{premise:696e73744e6f6e656d7074794f66496e68616269746564}{P32.}] \refstepcounter{paperpremise}\label{premise:696e73744e6f6e656d7074794f66496e68616269746564}\textbf{Inst nonempty of inhabited.} Let \(\alpha\) be a sort. Fix a structure \(\mathit{inst}\) of type \(\operatorname{Inhabited}\left(\alpha\right)\). Then \(\alpha\text{ is nonempty}\).
\item[\paperlineanchor{premise:6d756c5f6f6e65}\paperlink{premise:6d756c5f6f6e65}{P33.}] \refstepcounter{paperpremise}\label{premise:6d756c5f6f6e65}\textbf{Product one.} Let \(M\) be a type. Fix a structure \(\mathit{inst}\) of type \(\operatorname{Mul\,One\,Class}\left(M\right)\). Let \(a\) be an element of \(M\). Then \(a \cdot 1 = a\).
\item[\paperlineanchor{premise:747275655f616e64}\paperlink{premise:747275655f616e64}{P34.}] \refstepcounter{paperpremise}\label{premise:747275655f616e64}\textbf{True and.} Let \(p\) be a proposition. Then \(\left(\mathrm{True} \land p\right) = p\).
\end{description}

\section{Bipartite relations}
\begin{theorem}[{Incoming neighbors after reversing the relation}]\label{theorem:1}
Let \(\alpha\) and \(\beta\) be types. Let \(r\) be a relation from \(\alpha\) to \(\beta\). Let \(t\) be a finite subset of \(\beta\). Let \(a\) be an element of \(\alpha\). Assume the predicate \(r\left(a\right)\) is decidable. Then \(N^{-}_{\operatorname{swap}\left(r\right), t}(a) = N^{+}_{r, t}(a)\).
\end{theorem}

\begin{proof}
By definition, we obtain \(N^{-}_{\operatorname{swap}\left(r\right), t}(a) = N^{+}_{r, t}(a)\).
\end{proof}

\begin{theorem}[{Outgoing neighbors after reversing the relation}]\label{theorem:2}
Let \(\alpha\) and \(\beta\) be types. Let \(r\) be a relation from \(\alpha\) to \(\beta\). Let \(s\) be a finite subset of \(\alpha\). Let \(b\) be an element of \(\beta\). Fix a structure \(\mathit{inst}\) of type \(\Pi_{a\, :\, \alpha}\, \operatorname{Decidable}\left(r\left(a, b\right)\right)\). Then \(N^{+}_{\operatorname{swap}\left(r\right), s}(b) = N^{-}_{r, s}(b)\).
\end{theorem}

\begin{proof}
By definition, we obtain \(N^{+}_{\operatorname{swap}\left(r\right), s}(b) = N^{-}_{r, s}(b)\).
\end{proof}

\begin{theorem}[{The underlying set of incoming neighbors}]\label{theorem:3}
Let \(\alpha\) and \(\beta\) be types. Let \(r\) be a relation from \(\alpha\) to \(\beta\). Let \(s\) be a finite subset of \(\alpha\). Let \(b\) be an element of \(\beta\). Fix a structure \(\mathit{inst}\) of type \(\Pi_{a\, :\, \alpha}\, \operatorname{Decidable}\left(r\left(a, b\right)\right)\). Then
\begin{center}
\(\displaystyle \operatorname{underlying\,set}(N^{-}_{r,\allowbreak{} s}(b)) = \{a\in \alpha \mid a \in s \land r(a,\allowbreak{} b)\}.\)
\end{center}
\end{theorem}

\begin{proof}
By \paperlink{premise:46696e7365742e636f655f66696c746572}{Prerequisite~P\ref{premise:46696e7365742e636f655f66696c746572}} with \(p = a \mapsto r\left(a, b\right)\), we obtain \(\operatorname{underlying\,set}(N^{-}_{r, s}(b)) = \{a\in \alpha \mid a \in s \land r\left(a, b\right)\}\).
\end{proof}

\begin{lemma}[{The underlying set of outgoing neighbors}]\label{theorem:4}
Let \(\alpha\) and \(\beta\) be types. Let \(r\) be a relation from \(\alpha\) to \(\beta\). Let \(t\) be a finite subset of \(\beta\). Let \(a\) be an element of \(\alpha\). Assume the predicate \(r\left(a\right)\) is decidable. Then
\begin{center}
\(\displaystyle \operatorname{underlying\,set}(N^{+}_{r,\allowbreak{} t}(a)) = \{b\in \beta \mid b \in t \land r(a,\allowbreak{} b)\}.\)
\end{center}
\end{lemma}

\begin{proof}
By \paperlink{premise:46696e7365742e636f655f66696c746572}{Prerequisite~P\ref{premise:46696e7365742e636f655f66696c746572}} with \(p = b \mapsto r\left(a, b\right)\) and \(s = t\), we obtain \(\operatorname{underlying\,set}(N^{+}_{r, t}(a)) = \{b\in \beta \mid b \in t \land r\left(a, b\right)\}\).
\end{proof}

\begin{lemma}[{Membership in an incoming neighborhood}]\label{theorem:5}
Let \(\alpha\) and \(\beta\) be types. Let \(r\) be a relation from \(\alpha\) to \(\beta\). Let \(s\) be a finite subset of \(\alpha\). Let \(b\) be an element of \(\beta\). Fix a structure \(\mathit{inst}\) of type \(\Pi_{a\, :\, \alpha}\, \operatorname{Decidable}\left(r\left(a, b\right)\right)\). Let \(a\) be an element of \(\alpha\). Then \(a \in N^{-}_{r, s}(b) \iff a \in s \land r\left(a, b\right)\).
\end{lemma}

\begin{proof}
By \paperlink{premise:46696e7365742e6d656d5f66696c746572}{Prerequisite~P\ref{premise:46696e7365742e6d656d5f66696c746572}}, we obtain \(a \in N^{-}_{r, s}(b) \iff a \in s \land r\left(a, b\right)\).
\end{proof}

\begin{theorem}[{Membership in an outgoing neighborhood}]\label{theorem:6}
Let \(\alpha\) and \(\beta\) be types. Let \(r\) be a relation from \(\alpha\) to \(\beta\). Let \(t\) be a finite subset of \(\beta\). Let \(a\) be an element of \(\alpha\). Assume the predicate \(r\left(a\right)\) is decidable. Let \(b\) be an element of \(\beta\). Then \(b \in N^{+}_{r, t}(a) \iff b \in t \land r\left(a, b\right)\).
\end{theorem}

\begin{proof}
By \paperlink{premise:46696e7365742e6d656d5f66696c746572}{Prerequisite~P\ref{premise:46696e7365742e6d656d5f66696c746572}}, we obtain \(b \in N^{+}_{r, t}(a) \iff b \in t \land r\left(a, b\right)\).
\end{proof}

\begin{lemma}[{A sum over incidences in two orders}]\label{theorem:7}
Let \(R\), \(\alpha\), and \(\beta\) be types. Let \(r\) be a relation from \(\alpha\) to \(\beta\). Let \(s\) be a finite subset of \(\alpha\). Let \(t\) be a finite subset of \(\beta\). Equip \(R\) with a commutative additive monoid structure. Let \(f\) be a function of type \(\alpha \to \beta \to R\). Assume the relation \(r\) is decidable. Then
\begin{center}
\(\displaystyle \sum_{a\in s}(\sum_{b\in N^{+}_{r,\allowbreak{} t}(a)}(f(a,\allowbreak{} b))) = \sum_{b\in t}(\sum_{a\in N^{-}_{r,\allowbreak{} s}(b)}(f(a,\allowbreak{} b))).\)
\end{center}
\end{lemma}

\begin{proof}
The substitution uses the following argument.
\begin{proofstep}
Fix \(x\) in \(\alpha\).
Assume \(x \in s\).
By definition, we obtain \(\sum_{b\in \{b\in t \mid r(x,\allowbreak{} b)\}}(f(x,\allowbreak{} b)) = \sum_{b\in \{b\in t \mid r(x,\allowbreak{} b)\}}(f(x,\allowbreak{} b))\).
\end{proofstep}
The substitution uses the following argument.
\begin{proofstep}
Fix \(x\) in \(\beta\).
Assume \(x \in t\).
By definition, we obtain \(\sum_{a_{1}\in \{a_{1}\in s \mid r(a_{1},\allowbreak{} x)\}}(f(a_{1},\allowbreak{} x)) = \sum_{a_{1}\in \{a_{1}\in s \mid r(a_{1},\allowbreak{} x)\}}(f(a_{1},\allowbreak{} x))\).
\end{proofstep}
The substitution uses the following argument.
\begin{proofstep}
Fix \(x\) in \(\alpha\).
Assume \(x \in s\).
By \paperlink{premise:46696e7365742e73756d5f66696c746572}{Prerequisite~P\ref{premise:46696e7365742e73756d5f66696c746572}} with \(p = r\left(x\right)\) and \(f\left(x\right)\) in place of its parameter \(f\), we obtain \(\sum_{a_{1}\in \{a_{1}\in t \mid r(x,\allowbreak{} a_{1})\}}(f(x,\allowbreak{} a_{1})) = \sum_{a_{1}\in t}(\begin{cases}f(x,\allowbreak{} a_{1}) &\text{if } r(x,\allowbreak{} a_{1})\\ 0 &\text{otherwise}\end{cases})\).
\end{proofstep}
The substitution uses the following argument.
\begin{proofstep}
Fix \(x\) in \(\beta\).
Assume \(x \in t\).
By \paperlink{premise:46696e7365742e73756d5f66696c746572}{Prerequisite~P\ref{premise:46696e7365742e73756d5f66696c746572}} with \(p = a_{1} \mapsto r\left(a_{1}, x\right)\) and \(a_{1} \mapsto f\left(a_{1}, x\right)\) in place of its parameter \(f\), we obtain \(\sum_{a_{1}\in \{a_{1}\in s \mid r(a_{1},\allowbreak{} x)\}}(f(a_{1},\allowbreak{} x)) = \sum_{a_{1}\in s}(\begin{cases}f(a_{1},\allowbreak{} x) &\text{if } r(a_{1},\allowbreak{} x)\\ 0 &\text{otherwise}\end{cases})\).
\end{proofstep}
By \paperlink{premise:46696e7365742e73756d5f636f6d6d}{Prerequisite~P\ref{premise:46696e7365742e73756d5f636f6d6d}}, we obtain \(\sum_{x\in s}(\sum_{a\in t}(\begin{cases}f(x,\allowbreak{} a) &\text{if } r(x,\allowbreak{} a)\\ 0 &\text{otherwise}\end{cases})) = \sum_{x\in t}(\sum_{a\in s}(\begin{cases}f(a,\allowbreak{} x) &\text{if } r(a,\allowbreak{} x)\\ 0 &\text{otherwise}\end{cases}))\).
Rewriting with \paperlink{premise:46696e7365742e73756d5f636f6e6772}{Prerequisite~P\ref{premise:46696e7365742e73756d5f636f6e6772}} and the two identities established above in the preceding result gives \(\sum_{x\in s}(\sum_{b\in \{b\in t \mid r(x,\allowbreak{} b)\}}(f(x,\allowbreak{} b))) = \sum_{x\in t}(\sum_{a\in \{a\in s \mid r(a,\allowbreak{} x)\}}(f(a,\allowbreak{} x)))\).
Rewriting with \paperlink{premise:46696e7365742e73756d5f636f6e6772}{Prerequisite~P\ref{premise:46696e7365742e73756d5f636f6e6772}} and the identity established above in the preceding result gives \(\sum_{x\in s}(\sum_{b\in \{b\in t \mid r(x,\allowbreak{} b)\}}(f(x,\allowbreak{} b))) = \sum_{b\in t}(\sum_{a\in N^{-}_{r,\allowbreak{} s}(b)}(f(a,\allowbreak{} b)))\).
Rewriting with \paperlink{premise:46696e7365742e73756d5f636f6e6772}{Prerequisite~P\ref{premise:46696e7365742e73756d5f636f6e6772}} and the identity established above in the preceding result gives \(\sum_{a\in s}(\sum_{b\in N^{+}_{r,\allowbreak{} t}(a)}(f(a,\allowbreak{} b))) = \sum_{b\in t}(\sum_{a\in N^{-}_{r,\allowbreak{} s}(b)}(f(a,\allowbreak{} b)))\).
\end{proof}

\begin{theorem}[{A product over incidences in two orders}]\label{theorem:8}
Let \(R\), \(\alpha\), and \(\beta\) be types. Let \(r\) be a relation from \(\alpha\) to \(\beta\). Let \(s\) be a finite subset of \(\alpha\). Let \(t\) be a finite subset of \(\beta\). Fix a structure \(\mathit{inst}\) of type \(\operatorname{Comm\,Monoid}\left(R\right)\). Let \(f\) be a function of type \(\alpha \to \beta \to R\). Assume the relation \(r\) is decidable. Then
\begin{center}
\(\displaystyle \prod_{a\in s}(\prod_{b\in N^{+}_{r,\allowbreak{} t}(a)}(f(a,\allowbreak{} b))) = \prod_{b\in t}(\prod_{a\in N^{-}_{r,\allowbreak{} s}(b)}(f(a,\allowbreak{} b))).\)
\end{center}
\end{theorem}

\begin{proof}
The substitution uses the following argument.
\begin{proofstep}
Fix \(x\) in \(\alpha\).
Assume \(x \in s\).
By definition, we obtain \(\prod_{b\in \{b\in t \mid r(x,\allowbreak{} b)\}}(f(x,\allowbreak{} b)) = \prod_{b\in \{b\in t \mid r(x,\allowbreak{} b)\}}(f(x,\allowbreak{} b))\).
\end{proofstep}
The substitution uses the following argument.
\begin{proofstep}
Fix \(x\) in \(\beta\).
Assume \(x \in t\).
By definition, we obtain \(\prod_{a_{1}\in \{a_{1}\in s \mid r(a_{1},\allowbreak{} x)\}}(f(a_{1},\allowbreak{} x)) = \prod_{a_{1}\in \{a_{1}\in s \mid r(a_{1},\allowbreak{} x)\}}(f(a_{1},\allowbreak{} x))\).
\end{proofstep}
The substitution uses the following argument.
\begin{proofstep}
Fix \(x\) in \(\alpha\).
Assume \(x \in s\).
By \paperlink{premise:46696e7365742e70726f645f66696c746572}{Prerequisite~P\ref{premise:46696e7365742e70726f645f66696c746572}} with \(p = r\left(x\right)\) and \(f\left(x\right)\) in place of its parameter \(f\), we obtain \(\prod_{a_{1}\in \{a_{1}\in t \mid r(x,\allowbreak{} a_{1})\}}(f(x,\allowbreak{} a_{1})) = \prod_{a_{1}\in t}(\begin{cases}f(x,\allowbreak{} a_{1}) &\text{if } r(x,\allowbreak{} a_{1})\\ 1 &\text{otherwise}\end{cases})\).
\end{proofstep}
The substitution uses the following argument.
\begin{proofstep}
Fix \(x\) in \(\beta\).
Assume \(x \in t\).
By \paperlink{premise:46696e7365742e70726f645f66696c746572}{Prerequisite~P\ref{premise:46696e7365742e70726f645f66696c746572}} with \(p = a_{1} \mapsto r\left(a_{1}, x\right)\) and \(a_{1} \mapsto f\left(a_{1}, x\right)\) in place of its parameter \(f\), we obtain \(\prod_{a_{1}\in \{a_{1}\in s \mid r(a_{1},\allowbreak{} x)\}}(f(a_{1},\allowbreak{} x)) = \prod_{a_{1}\in s}(\begin{cases}f(a_{1},\allowbreak{} x) &\text{if } r(a_{1},\allowbreak{} x)\\ 1 &\text{otherwise}\end{cases})\).
\end{proofstep}
By \paperlink{premise:46696e7365742e70726f645f636f6d6d}{Prerequisite~P\ref{premise:46696e7365742e70726f645f636f6d6d}}, we obtain \(\prod_{x\in s}(\prod_{a\in t}(\begin{cases}f(x,\allowbreak{} a) &\text{if } r(x,\allowbreak{} a)\\ 1 &\text{otherwise}\end{cases})) = \prod_{x\in t}(\prod_{a\in s}(\begin{cases}f(a,\allowbreak{} x) &\text{if } r(a,\allowbreak{} x)\\ 1 &\text{otherwise}\end{cases}))\).
Rewriting with \paperlink{premise:46696e7365742e70726f645f636f6e6772}{Prerequisite~P\ref{premise:46696e7365742e70726f645f636f6e6772}} and the two identities established above in the preceding result gives \(\prod_{x\in s}(\prod_{b\in \{b\in t \mid r(x,\allowbreak{} b)\}}(f(x,\allowbreak{} b))) = \prod_{x\in t}(\prod_{a\in \{a\in s \mid r(a,\allowbreak{} x)\}}(f(a,\allowbreak{} x)))\).
Rewriting with \paperlink{premise:46696e7365742e70726f645f636f6e6772}{Prerequisite~P\ref{premise:46696e7365742e70726f645f636f6e6772}} and the identity established above in the preceding result gives \(\prod_{x\in s}(\prod_{b\in \{b\in t \mid r(x,\allowbreak{} b)\}}(f(x,\allowbreak{} b))) = \prod_{b\in t}(\prod_{a\in N^{-}_{r,\allowbreak{} s}(b)}(f(a,\allowbreak{} b)))\).
Rewriting with \paperlink{premise:46696e7365742e70726f645f636f6e6772}{Prerequisite~P\ref{premise:46696e7365742e70726f645f636f6e6772}} and the identity established above in the preceding result gives \(\prod_{a\in s}(\prod_{b\in N^{+}_{r,\allowbreak{} t}(a)}(f(a,\allowbreak{} b))) = \prod_{b\in t}(\prod_{a\in N^{-}_{r,\allowbreak{} s}(b)}(f(a,\allowbreak{} b)))\).
\end{proof}

\begin{lemma}[{Counting incidences by rows and columns}]\label{theorem:9}
Let \(\alpha\) and \(\beta\) be types. Let \(r\) be a relation from \(\alpha\) to \(\beta\). Let \(s\) be a finite subset of \(\alpha\). Let \(t\) be a finite subset of \(\beta\). Assume the relation \(r\) is decidable. Then
\begin{center}
\(\displaystyle \sum_{a\in s}(|N^{+}_{r,\allowbreak{} t}(a)|) = \sum_{b\in t}(|N^{-}_{r,\allowbreak{} s}(b)|).\)
\end{center}
\end{lemma}

\begin{proof}
The substitution uses the following argument.
\begin{proofstep}
Fix \(x\) in \(\alpha\).
Assume \(x \in s\).
By \paperlink{premise:46696e7365742e636172645f65715f73756d5f6f6e6573}{Prerequisite~P\ref{premise:46696e7365742e636172645f65715f73756d5f6f6e6573}} with \(N^{+}_{r, t}(x)\) in place of its parameter \(s\), we obtain \(\left|N^{+}_{r, t}(x)\right| = \sum_{x_{1}\in N^{+}_{r, t}(x)}\left(1\right)\).
\end{proofstep}
The substitution uses the following argument.
\begin{proofstep}
Fix \(x\) in \(\beta\).
Assume \(x \in t\).
By \paperlink{premise:46696e7365742e636172645f65715f73756d5f6f6e6573}{Prerequisite~P\ref{premise:46696e7365742e636172645f65715f73756d5f6f6e6573}} with \(N^{-}_{r, s}(x)\) in place of its parameter \(s\), we obtain \(\left|N^{-}_{r, s}(x)\right| = \sum_{x_{1}\in N^{-}_{r, s}(x)}\left(1\right)\).
\end{proofstep}
Rewriting with \paperlink{theorem:7}{Lemma~\ref{theorem:7}} gives \(\sum_{x\in s}(\sum_{x_{1}\in N^{+}_{r,\allowbreak{} t}(x)}(1)) = \sum_{x\in t}(\sum_{x_{1}\in N^{-}_{r,\allowbreak{} s}(x)}(1))\).
Rewriting with \paperlink{premise:46696e7365742e73756d5f636f6e6772}{Prerequisite~P\ref{premise:46696e7365742e73756d5f636f6e6772}} and the two identities established above in the preceding result gives \(\sum_{a\in s}(|N^{+}_{r,\allowbreak{} t}(a)|) = \sum_{b\in t}(|N^{-}_{r,\allowbreak{} s}(b)|)\).
\end{proof}

\subsection{Bounds in ordered semirings}
\begin{lemma}[{Double counting with a scalar lower bound on row degrees}]\label{theorem:10}
Let \(R\), \(\alpha\), and \(\beta\) be types. Let \(r\) be a relation from \(\alpha\) to \(\beta\). Let \(s\) be a finite subset of \(\alpha\). Let \(t\) be a finite subset of \(\beta\). Equip \(R\) with a semiring structure. Equip \(R\) with a partial order. Let \(\mathit{inst}_{2}\) be a proof of \(\operatorname{Is\,Ordered\,Ring}\left(R\right)\). Let \(m\) and \(n\) be elements of \(R\). Assume the relation \(r\) is decidable. Assume \(\forall a\, :\, \alpha,\; a \in s \implies m \le \operatorname{cast}_{R}(|N^{+}_{r,\allowbreak{} t}(a)|)\). Assume \(\forall b\, :\, \beta,\; b \in t \implies \operatorname{cast}_{R}(|N^{-}_{r,\allowbreak{} s}(b)|) \le n\). Then \(\left|s\right| \mathbin{\bullet} m \le \left|t\right| \mathbin{\bullet} n\).
\end{lemma}

\begin{proof}
Use the hypothesis \(\operatorname{Is\,Ordered\,Ring}\left(R\right)\).
By \paperlink{premise:49734f72646572656452696e672e746f49734f7264657265644164644d6f6e6f6964}{Prerequisite~P\ref{premise:49734f72646572656452696e672e746f49734f7264657265644164644d6f6e6f6964}}, we obtain \(\text{translation preserves order on }R\).
Use the hypothesis \(\forall a\, :\, \alpha,\; a \in s \implies m \le \operatorname{cast}_{R}(|N^{+}_{r,\allowbreak{} t}(a)|)\).
By \paperlink{premise:46696e7365742e636172645f6e736d756c5f6c655f73756d}{Prerequisite~P\ref{premise:46696e7365742e636172645f6e736d756c5f6c655f73756d}} with \(f = a \mapsto \operatorname{cast}_{R}\left(\left|N^{+}_{r, t}(a)\right|\right)\) and \(m\) in place of its parameter \(n\), we obtain \(|s| \mathbin{\bullet} m \le \sum_{a\in s}(\operatorname{cast}_{R}(|N^{+}_{r,\allowbreak{} t}(a)|))\).
Rewriting with \paperlink{theorem:9}{Lemma~\ref{theorem:9}} and \paperlink{premise:4e61742e636173745f73756d}{Prerequisite~P\ref{premise:4e61742e636173745f73756d}} gives \(\sum_{a\in s}(\operatorname{cast}_{R}(|N^{+}_{r,\allowbreak{} t}(a)|)) = \sum_{b\in t}(\operatorname{cast}_{R}(|N^{-}_{r,\allowbreak{} s}(b)|))\).
By transitivity of the displayed relations, we obtain \(|s| \mathbin{\bullet} m \le \sum_{b\in t}(\operatorname{cast}_{R}(|N^{-}_{r,\allowbreak{} s}(b)|))\).
Use the hypothesis \(\operatorname{Is\,Ordered\,Ring}\left(R\right)\).
By \paperlink{premise:49734f72646572656452696e672e746f49734f7264657265644164644d6f6e6f6964}{Prerequisite~P\ref{premise:49734f72646572656452696e672e746f49734f7264657265644164644d6f6e6f6964}}, we obtain \(\text{translation preserves order on }R\).
Use the hypothesis \(\forall b\, :\, \beta,\; b \in t \implies \operatorname{cast}_{R}(|N^{-}_{r,\allowbreak{} s}(b)|) \le n\).
By \paperlink{premise:46696e7365742e73756d5f6c655f636172645f6e736d756c}{Prerequisite~P\ref{premise:46696e7365742e73756d5f6c655f636172645f6e736d756c}} with \(t\) in place of its parameter \(s\) and \(f = b \mapsto \operatorname{cast}_{R}\left(\left|N^{-}_{r, s}(b)\right|\right)\), we obtain \(\sum_{b\in t}(\operatorname{cast}_{R}(|N^{-}_{r,\allowbreak{} s}(b)|)) \le |t| \mathbin{\bullet} n\).
By transitivity of the displayed relations, we obtain \(\left|s\right| \mathbin{\bullet} m \le \left|t\right| \mathbin{\bullet} n\).
\end{proof}

\begin{corollary}[{Double counting with a scalar lower bound on column degrees}]\label{theorem:11}
Let \(R\), \(\alpha\), and \(\beta\) be types. Let \(r\) be a relation from \(\alpha\) to \(\beta\). Let \(s\) be a finite subset of \(\alpha\). Let \(t\) be a finite subset of \(\beta\). Equip \(R\) with a semiring structure. Equip \(R\) with a partial order. Let \(\mathit{inst}_{2}\) be a proof of \(\operatorname{Is\,Ordered\,Ring}\left(R\right)\). Let \(m\) and \(n\) be elements of \(R\). Assume the relation \(r\) is decidable. Assume \(\forall b\, :\, \beta,\; b \in t \implies n \le \operatorname{cast}_{R}(|N^{-}_{r,\allowbreak{} s}(b)|)\). Assume \(\forall a\, :\, \alpha,\; a \in s \implies \operatorname{cast}_{R}(|N^{+}_{r,\allowbreak{} t}(a)|) \le m\). Then \(\left|t\right| \mathbin{\bullet} n \le \left|s\right| \mathbin{\bullet} m\).
\end{corollary}

\begin{proof}
Apply \paperlink{theorem:10}{Lemma~\ref{theorem:10}} with \(\operatorname{swap}\left(r\right)\) in place of its parameter \(r\).
\end{proof}

\subsection{Strict bounds in ordered semirings}
\begin{lemma}[{Double counting with a strict lower bound on row degrees}]\label{theorem:12}
Let \(R\), \(\alpha\), and \(\beta\) be types. Equip \(R\) with a semiring structure. Equip \(R\) with a partial order. Assume \(R\text{ is a strictly ordered semiring}\). Let \(r\) be a relation from \(\alpha\) to \(\beta\). Let \(s\) be a finite subset of \(\alpha\). Let \(t\) be a finite subset of \(\beta\). Let \(m\) and \(n\) be elements of \(R\). Assume the relation \(r\) is decidable. Assume \(s\ne\varnothing\). Assume \(\forall a\, :\, \alpha,\; a \in s \implies m < \operatorname{cast}_{R}(|N^{+}_{r,\allowbreak{} t}(a)|)\). Assume \(\forall b\, :\, \beta,\; b \in t \implies \operatorname{cast}_{R}(|N^{-}_{r,\allowbreak{} s}(b)|) \le n\). Then \(\left|s\right| \mathbin{\bullet} m < \left|t\right| \mathbin{\bullet} n\).
\end{lemma}

\begin{proof}
By definition, we obtain \(\left|s\right| \mathbin{\bullet} m = \left|s\right| \mathbin{\bullet} m\).
Rewriting with \paperlink{premise:46696e7365742e73756d5f636f6e7374}{Prerequisite~P\ref{premise:46696e7365742e73756d5f636f6e7374}} in the preceding result gives \(\left|s\right| \mathbin{\bullet} m = \sum_{x\in s}\left(m\right)\).
Use the hypothesis \(R\text{ is a strictly ordered semiring}\).
By \paperlink{premise:49735374726963744f72646572656452696e672e746f49734f72646572656443616e63656c4164644d6f6e6f6964}{Prerequisite~P\ref{premise:49735374726963744f72646572656452696e672e746f49734f72646572656443616e63656c4164644d6f6e6f6964}}, we obtain \(\text{translation preserves and reflects order on }R\).
Use the hypothesis \(s\ne\varnothing\).
Use the hypothesis \(\forall a\, :\, \alpha,\; a \in s \implies m < \operatorname{cast}_{R}(|N^{+}_{r,\allowbreak{} t}(a)|)\).
By \paperlink{premise:46696e7365742e73756d5f6c745f73756d5f6f665f6e6f6e656d707479}{Prerequisite~P\ref{premise:46696e7365742e73756d5f6c745f73756d5f6f665f6e6f6e656d707479}}, we obtain \(\sum_{i\in s}(m) < \sum_{i\in s}(\operatorname{cast}_{R}(|N^{+}_{r,\allowbreak{} t}(i)|))\).
By transitivity of the displayed relations, we obtain \(|s| \mathbin{\bullet} m < \sum_{a\in s}(\operatorname{cast}_{R}(|N^{+}_{r,\allowbreak{} t}(a)|))\).
Rewriting with \paperlink{theorem:9}{Lemma~\ref{theorem:9}}, \paperlink{premise:4e61742e636173745f73756d}{Prerequisite~P\ref{premise:4e61742e636173745f73756d}}, \paperlink{premise:4e61742e636173745f696e6a}{Prerequisite~P\ref{premise:4e61742e636173745f696e6a}} and \paperlink{premise:49735374726963744f72646572656452696e672e746f436861725a65726f}{Prerequisite~P\ref{premise:49735374726963744f72646572656452696e672e746f436861725a65726f}} (using \(R\text{ is a strictly ordered semiring}\)) gives \(\sum_{a\in s}(\operatorname{cast}_{R}(|N^{+}_{r,\allowbreak{} t}(a)|)) = \sum_{b\in t}(\operatorname{cast}_{R}(|N^{-}_{r,\allowbreak{} s}(b)|))\).
By transitivity of the displayed relations, we obtain \(|s| \mathbin{\bullet} m < \sum_{b\in t}(\operatorname{cast}_{R}(|N^{-}_{r,\allowbreak{} s}(b)|))\).
Use the hypothesis \(R\text{ is a strictly ordered semiring}\).
By \paperlink{premise:49735374726963744f72646572656452696e672e746f49734f72646572656452696e67}{Prerequisite~P\ref{premise:49735374726963744f72646572656452696e672e746f49734f72646572656452696e67}}, we obtain \(\operatorname{Is\,Ordered\,Ring}\left(R\right)\).
By \paperlink{premise:49734f72646572656452696e672e746f49734f7264657265644164644d6f6e6f6964}{Prerequisite~P\ref{premise:49734f72646572656452696e672e746f49734f7264657265644164644d6f6e6f6964}}, we obtain \(\text{translation preserves order on }R\).
Use the hypothesis \(\forall b\, :\, \beta,\; b \in t \implies \operatorname{cast}_{R}(|N^{-}_{r,\allowbreak{} s}(b)|) \le n\).
By \paperlink{premise:46696e7365742e73756d5f6c655f636172645f6e736d756c}{Prerequisite~P\ref{premise:46696e7365742e73756d5f6c655f636172645f6e736d756c}} with \(t\) in place of its parameter \(s\) and \(f = b \mapsto \operatorname{cast}_{R}\left(\left|N^{-}_{r, s}(b)\right|\right)\), we obtain \(\sum_{b\in t}(\operatorname{cast}_{R}(|N^{-}_{r,\allowbreak{} s}(b)|)) \le |t| \mathbin{\bullet} n\).
By transitivity of the displayed relations, we obtain \(\left|s\right| \mathbin{\bullet} m < \left|t\right| \mathbin{\bullet} n\).
\end{proof}

\begin{lemma}[{Double counting with a strict upper bound on column degrees}]\label{theorem:13}
Let \(R\), \(\alpha\), and \(\beta\) be types. Equip \(R\) with a semiring structure. Equip \(R\) with a partial order. Assume \(R\text{ is a strictly ordered semiring}\). Let \(r\) be a relation from \(\alpha\) to \(\beta\). Let \(s\) be a finite subset of \(\alpha\). Let \(t\) be a finite subset of \(\beta\). Let \(m\) and \(n\) be elements of \(R\). Assume the relation \(r\) is decidable. Assume \(t\ne\varnothing\). Assume \(\forall a\, :\, \alpha,\; a \in s \implies m \le \operatorname{cast}_{R}(|N^{+}_{r,\allowbreak{} t}(a)|)\). Assume \(\forall b\, :\, \beta,\; b \in t \implies \operatorname{cast}_{R}\left(\left|N^{-}_{r, s}(b)\right|\right) < n\). Then \(\left|s\right| \mathbin{\bullet} m < \left|t\right| \mathbin{\bullet} n\).
\end{lemma}

\begin{proof}
Use the hypothesis \(R\text{ is a strictly ordered semiring}\).
By \paperlink{premise:49735374726963744f72646572656452696e672e746f49734f72646572656452696e67}{Prerequisite~P\ref{premise:49735374726963744f72646572656452696e672e746f49734f72646572656452696e67}}, we obtain \(\operatorname{Is\,Ordered\,Ring}\left(R\right)\).
By \paperlink{premise:49734f72646572656452696e672e746f49734f7264657265644164644d6f6e6f6964}{Prerequisite~P\ref{premise:49734f72646572656452696e672e746f49734f7264657265644164644d6f6e6f6964}}, we obtain \(\text{translation preserves order on }R\).
Use the hypothesis \(\forall a\, :\, \alpha,\; a \in s \implies m \le \operatorname{cast}_{R}(|N^{+}_{r,\allowbreak{} t}(a)|)\).
By \paperlink{premise:46696e7365742e636172645f6e736d756c5f6c655f73756d}{Prerequisite~P\ref{premise:46696e7365742e636172645f6e736d756c5f6c655f73756d}} with \(f = a \mapsto \operatorname{cast}_{R}\left(\left|N^{+}_{r, t}(a)\right|\right)\) and \(m\) in place of its parameter \(n\), we obtain \(|s| \mathbin{\bullet} m \le \sum_{a\in s}(\operatorname{cast}_{R}(|N^{+}_{r,\allowbreak{} t}(a)|))\).
Rewriting with \paperlink{theorem:9}{Lemma~\ref{theorem:9}}, \paperlink{premise:4e61742e636173745f73756d}{Prerequisite~P\ref{premise:4e61742e636173745f73756d}}, \paperlink{premise:4e61742e636173745f696e6a}{Prerequisite~P\ref{premise:4e61742e636173745f696e6a}} and \paperlink{premise:49735374726963744f72646572656452696e672e746f436861725a65726f}{Prerequisite~P\ref{premise:49735374726963744f72646572656452696e672e746f436861725a65726f}} (using \(R\text{ is a strictly ordered semiring}\)) gives \(\sum_{a\in s}(\operatorname{cast}_{R}(|N^{+}_{r,\allowbreak{} t}(a)|)) = \sum_{b\in t}(\operatorname{cast}_{R}(|N^{-}_{r,\allowbreak{} s}(b)|))\).
By transitivity of the displayed relations, we obtain \(|s| \mathbin{\bullet} m \le \sum_{b\in t}(\operatorname{cast}_{R}(|N^{-}_{r,\allowbreak{} s}(b)|))\).
Use the hypothesis \(R\text{ is a strictly ordered semiring}\).
By \paperlink{premise:49735374726963744f72646572656452696e672e746f49734f72646572656443616e63656c4164644d6f6e6f6964}{Prerequisite~P\ref{premise:49735374726963744f72646572656452696e672e746f49734f72646572656443616e63656c4164644d6f6e6f6964}}, we obtain \(\text{translation preserves and reflects order on }R\).
Use the hypothesis \(t\ne\varnothing\).
Use the hypothesis \(\forall b\, :\, \beta,\; b \in t \implies \operatorname{cast}_{R}\left(\left|N^{-}_{r, s}(b)\right|\right) < n\).
By \paperlink{premise:46696e7365742e73756d5f6c745f73756d5f6f665f6e6f6e656d707479}{Prerequisite~P\ref{premise:46696e7365742e73756d5f6c745f73756d5f6f665f6e6f6e656d707479}}, we obtain \(\sum_{i\in t}(\operatorname{cast}_{R}(|N^{-}_{r,\allowbreak{} s}(i)|)) < \sum_{i\in t}(n)\).
By transitivity of the displayed relations, we obtain \(\left|s\right| \mathbin{\bullet} m < \sum_{x\in t}\left(n\right)\).
By \paperlink{premise:46696e7365742e73756d5f636f6e7374}{Prerequisite~P\ref{premise:46696e7365742e73756d5f636f6e7374}} with \(b = n\), we obtain \(\sum_{x\in t}\left(n\right) = \left|t\right| \mathbin{\bullet} n\).
By transitivity of the displayed relations, we obtain \(\left|s\right| \mathbin{\bullet} m < \left|t\right| \mathbin{\bullet} n\).
\end{proof}

\begin{corollary}[{Double counting with a strict lower bound on column degrees}]\label{theorem:14}
Let \(R\), \(\alpha\), and \(\beta\) be types. Equip \(R\) with a semiring structure. Equip \(R\) with a partial order. Assume \(R\text{ is a strictly ordered semiring}\). Let \(r\) be a relation from \(\alpha\) to \(\beta\). Let \(s\) be a finite subset of \(\alpha\). Let \(t\) be a finite subset of \(\beta\). Let \(m\) and \(n\) be elements of \(R\). Assume the relation \(r\) is decidable. Assume \(t\ne\varnothing\). Assume \(\forall b\, :\, \beta,\; b \in t \implies n < \operatorname{cast}_{R}\left(\left|N^{-}_{r, s}(b)\right|\right)\). Assume \(\forall a\, :\, \alpha,\; a \in s \implies \operatorname{cast}_{R}(|N^{+}_{r,\allowbreak{} t}(a)|) \le m\). Then \(\left|t\right| \mathbin{\bullet} n < \left|s\right| \mathbin{\bullet} m\).
\end{corollary}

\begin{proof}
Apply \paperlink{theorem:12}{Lemma~\ref{theorem:12}} with \(\operatorname{swap}\left(r\right)\) in place of its parameter \(r\).
\end{proof}

\begin{corollary}[{Double counting with a strict upper bound on row degrees}]\label{theorem:15}
Let \(R\), \(\alpha\), and \(\beta\) be types. Equip \(R\) with a semiring structure. Equip \(R\) with a partial order. Assume \(R\text{ is a strictly ordered semiring}\). Let \(r\) be a relation from \(\alpha\) to \(\beta\). Let \(s\) be a finite subset of \(\alpha\). Let \(t\) be a finite subset of \(\beta\). Let \(m\) and \(n\) be elements of \(R\). Assume the relation \(r\) is decidable. Assume \(s\ne\varnothing\). Assume \(\forall b\, :\, \beta,\; b \in t \implies n \le \operatorname{cast}_{R}(|N^{-}_{r,\allowbreak{} s}(b)|)\). Assume \(\forall a\, :\, \alpha,\; a \in s \implies \operatorname{cast}_{R}(|N^{+}_{r,\allowbreak{} t}(a)|) < m\). Then \(\left|t\right| \mathbin{\bullet} n < \left|s\right| \mathbin{\bullet} m\).
\end{corollary}

\begin{proof}
Apply \paperlink{theorem:13}{Lemma~\ref{theorem:13}} with \(\operatorname{swap}\left(r\right)\) in place of its parameter \(r\).
\end{proof}

\subsection{Further results}
\begin{corollary}[{Double counting with a natural-number lower bound on row degrees}]\label{theorem:16}
Let \(\alpha\) and \(\beta\) be types. Let \(r\) be a relation from \(\alpha\) to \(\beta\). Let \(s\) be a finite subset of \(\alpha\). Let \(t\) be a finite subset of \(\beta\). Let \(m\) and \(n\) be natural numbers. Assume the relation \(r\) is decidable. Assume \(\forall a\, :\, \alpha,\; a \in s \implies m \le \left|N^{+}_{r, t}(a)\right|\). Assume \(\forall b\, :\, \beta,\; b \in t \implies \left|N^{-}_{r, s}(b)\right| \le n\). Then \(\left|s\right| \cdot m \le \left|t\right| \cdot n\).
\end{corollary}

\begin{proof}
By \paperlink{premise:4e61742e696e737449735374726963744f72646572656452696e67}{Prerequisite~P\ref{premise:4e61742e696e737449735374726963744f72646572656452696e67}}, we obtain \(\mathbb{N}\text{ is a strictly ordered semiring}\).
By \paperlink{premise:49735374726963744f72646572656452696e672e746f49734f72646572656452696e67}{Prerequisite~P\ref{premise:49735374726963744f72646572656452696e672e746f49734f72646572656452696e67}} with \(R = \mathbb{N}\), we obtain \(\operatorname{Is\,Ordered\,Ring}\left(\mathbb{N}\right)\).
Apply \paperlink{theorem:10}{Lemma~\ref{theorem:10}} with \(R = \mathbb{N}\).
\end{proof}

\begin{corollary}[{Double counting with a natural-number lower bound on column degrees}]\label{theorem:17}
Let \(\alpha\) and \(\beta\) be types. Let \(r\) be a relation from \(\alpha\) to \(\beta\). Let \(s\) be a finite subset of \(\alpha\). Let \(t\) be a finite subset of \(\beta\). Let \(m\) and \(n\) be natural numbers. Assume the relation \(r\) is decidable. Assume \(\forall b\, :\, \beta,\; b \in t \implies n \le \left|N^{-}_{r, s}(b)\right|\). Assume \(\forall a\, :\, \alpha,\; a \in s \implies \left|N^{+}_{r, t}(a)\right| \le m\). Then \(\left|t\right| \cdot n \le \left|s\right| \cdot m\).
\end{corollary}

\begin{proof}
By \paperlink{premise:4e61742e696e737449735374726963744f72646572656452696e67}{Prerequisite~P\ref{premise:4e61742e696e737449735374726963744f72646572656452696e67}}, we obtain \(\mathbb{N}\text{ is a strictly ordered semiring}\).
By \paperlink{premise:49735374726963744f72646572656452696e672e746f49734f72646572656452696e67}{Prerequisite~P\ref{premise:49735374726963744f72646572656452696e672e746f49734f72646572656452696e67}} with \(R = \mathbb{N}\), we obtain \(\operatorname{Is\,Ordered\,Ring}\left(\mathbb{N}\right)\).
Apply \paperlink{theorem:11}{Corollary~\ref{theorem:11}} with \(R = \mathbb{N}\).
\end{proof}

\begin{theorem}[{An equality of weighted cardinalities}]\label{theorem:18}
Let \(\alpha\) and \(\beta\) be types. Let \(r\) be a relation from \(\alpha\) to \(\beta\). Let \(s\) be a finite subset of \(\alpha\). Let \(t\) be a finite subset of \(\beta\). Let \(m\) and \(n\) be natural numbers. Assume the relation \(r\) is decidable. Assume \(\forall a\, :\, \alpha,\; a \in s \implies \left|N^{+}_{r, t}(a)\right| = m\). Assume \(\forall b\, :\, \beta,\; b \in t \implies \left|N^{-}_{r, s}(b)\right| = n\). Then \(\left|s\right| \cdot m = \left|t\right| \cdot n\).
\end{theorem}

\begin{proof}
For the required implication, argue as follows.
\begin{proofstep}
Fix \(a\) in \(\alpha\).
Assume \(a \in s\).
Use the hypothesis \(a \in s\).
Applying the hypothesis \(\forall a\, :\, \alpha,\; a \in s \implies \left|N^{+}_{r, t}(a)\right| = m\), we obtain \(\left|N^{+}_{r, t}(a)\right| = m\).
By \paperlink{premise:45712e6765}{Prerequisite~P\ref{premise:45712e6765}} with \(\mathbb{N}\) as its type parameter \(\alpha\), we obtain \(m \le \left|N^{+}_{r, t}(a)\right|\).
\end{proofstep}
For the required implication, argue as follows.
\begin{proofstep}
Fix \(b\) in \(\beta\).
Assume \(b \in t\).
Use the hypothesis \(b \in t\).
Applying the hypothesis \(\forall b\, :\, \beta,\; b \in t \implies \left|N^{-}_{r, s}(b)\right| = n\), we obtain \(\left|N^{-}_{r, s}(b)\right| = n\).
By \paperlink{premise:45712e6c65}{\emph{Alias}} (\paperlink{premise:45712e6c65}{Prerequisite~P\ref{premise:45712e6c65}}) with \(\mathbb{N}\) as its type parameter \(\alpha\), we obtain \(\left|N^{-}_{r, s}(b)\right| \le n\).
\end{proofstep}
By \paperlink{theorem:16}{Corollary~\ref{theorem:16}}, we obtain \(\left|s\right| \cdot m \le \left|t\right| \cdot n\).
For the required implication, argue as follows.
\begin{proofstep}
Fix \(a\) in \(\beta\).
Assume \(a \in t\).
Use the hypothesis \(a \in t\).
Applying the hypothesis \(\forall b\, :\, \beta,\; b \in t \implies \left|N^{-}_{r, s}(b)\right| = n\), we obtain \(\left|N^{-}_{r, s}(a)\right| = n\).
By \paperlink{premise:45712e6765}{Prerequisite~P\ref{premise:45712e6765}} with \(\mathbb{N}\) as its type parameter \(\alpha\), we obtain \(n \le \left|N^{-}_{r, s}(a)\right|\).
\end{proofstep}
For the required implication, argue as follows.
\begin{proofstep}
Fix \(b\) in \(\alpha\).
Assume \(b \in s\).
Use the hypothesis \(b \in s\).
Applying the hypothesis \(\forall a\, :\, \alpha,\; a \in s \implies \left|N^{+}_{r, t}(a)\right| = m\), we obtain \(\left|N^{+}_{r, t}(b)\right| = m\).
By \paperlink{premise:45712e6c65}{\emph{Alias}} (\paperlink{premise:45712e6c65}{Prerequisite~P\ref{premise:45712e6c65}}) with \(\mathbb{N}\) as its type parameter \(\alpha\), we obtain \(\left|N^{+}_{r, t}(b)\right| \le m\).
\end{proofstep}
By \paperlink{theorem:17}{Corollary~\ref{theorem:17}}, we obtain \(\left|t\right| \cdot n \le \left|s\right| \cdot m\).
By \paperlink{premise:4c452e6c652e616e746973796d6d}{\emph{Alias}} (\paperlink{premise:4c452e6c652e616e746973796d6d}{Prerequisite~P\ref{premise:4c452e6c652e616e746973796d6d}}) with \(\mathbb{N}\) as its type parameter \(\alpha\), we obtain \(\left|s\right| \cdot m = \left|t\right| \cdot n\).
\end{proof}

\begin{lemma}[{Bounding nonempty rows by columns of degree at most one}]\label{theorem:19}
Let \(\alpha\) and \(\beta\) be types. Let \(r\) be a relation from \(\alpha\) to \(\beta\). Let \(s\) be a finite subset of \(\alpha\). Let \(t\) be a finite subset of \(\beta\). Assume \(\forall a\, :\, \alpha,\; a \in s \implies (\exists b\, :\, \beta,\; b \in t \land r(a,\allowbreak{} b))\). Assume \(\forall b\, :\, \beta,\; b \in t \implies \operatorname{Subsingleton}(\{a_{1}\in \alpha \mid a_{1} \in s \land r(a_{1},\allowbreak{} b)\})\). Then \(\left|s\right| \le \left|t\right|\).
\end{lemma}

\begin{proof}
For the required implication, argue as follows.
\begin{proofstep}
Fix \(a\) in \(\alpha\).
Assume \(a \in s\).
By \paperlink{premise:46696e7365742e636172645f706f73}{Prerequisite~P\ref{premise:46696e7365742e636172645f706f73}}, we obtain \(0 < \left|N^{+}_{r, t}(a)\right| \iff N^{+}_{r, t}(a)\ne\varnothing\).
Use the hypothesis \(a \in s\).
Applying the hypothesis \(\forall a\, :\, \alpha,\; a \in s \implies (\exists b\, :\, \beta,\; b \in t \land r(a,\allowbreak{} b))\), we obtain \(\{b\in \beta \mid b \in t \land r\left(a, b\right)\}\ne\varnothing\).
Rewriting with \paperlink{theorem:4}{Lemma~\ref{theorem:4}} in the preceding result gives \(\operatorname{underlying\,set}(N^{+}_{r, t}(a))\ne\varnothing\).
Rewriting with \paperlink{premise:46696e7365742e636f655f6e6f6e656d707479}{Prerequisite~P\ref{premise:46696e7365742e636f655f6e6f6e656d707479}} in the preceding result gives \(N^{+}_{r, t}(a)\ne\varnothing\).
By the reverse implication, we obtain \(0 < \left|N^{+}_{r, t}(a)\right|\).
\end{proofstep}
For the required implication, argue as follows.
\begin{proofstep}
Fix \(b\) in \(\beta\).
Assume \(b \in t\).
By \paperlink{premise:46696e7365742e636172645f6c655f6f6e65}{Prerequisite~P\ref{premise:46696e7365742e636172645f6c655f6f6e65}}, we obtain \(|N^{-}_{r,\allowbreak{} s}(b)| \le 1 \iff (\forall a\, :\, \alpha,\; a \in N^{-}_{r,\allowbreak{} s}(b) \implies (\forall b_{1}\, :\, \alpha,\; b_{1} \in N^{-}_{r,\allowbreak{} s}(b) \implies a = b_{1}))\).
The substitution uses the following argument.
\begin{proofstep}
Fix \(a\) in \(\alpha\).
By \paperlink{theorem:5}{Lemma~\ref{theorem:5}}, we obtain \(a \in N^{-}_{r, s}(b) \iff a \in s \land r\left(a, b\right)\).
By propositional extensionality, we obtain \(\left(a \in N^{-}_{r, s}(b)\right) = \left(a \in s \land r\left(a, b\right)\right)\).
For the required implication, argue as follows.
\begin{proofstep}
Fix \(b_{1}\) in \(\alpha\).
By \paperlink{theorem:5}{Lemma~\ref{theorem:5}}, we obtain \(b_{1} \in N^{-}_{r, s}(b) \iff b_{1} \in s \land r\left(b_{1}, b\right)\).
By propositional extensionality, we obtain \(\left(b_{1} \in N^{-}_{r, s}(b)\right) = \left(b_{1} \in s \land r\left(b_{1}, b\right)\right)\).
By definition, we obtain \(\left(a = b_{1}\right) = \left(a = b_{1}\right)\).
By \paperlink{premise:696d706c6965735f636f6e6772}{Prerequisite~P\ref{premise:696d706c6965735f636f6e6772}} with \(p_{1} = b_{1} \in N^{-}_{r, s}(b)\) and \(p_{2} = b_{1} \in s \land r\left(b_{1}, b\right)\) and \(q_{1} = a = b_{1}\) and \(q_{2} = a = b_{1}\), we obtain \((b_{1} \in N^{-}_{r,\allowbreak{} s}(b) \implies a = b_{1}) = (b_{1} \in s \land r(b_{1},\allowbreak{} b) \implies a = b_{1})\).
\end{proofstep}
By \paperlink{premise:666f72616c6c5f636f6e6772}{Prerequisite~P\ref{premise:666f72616c6c5f636f6e6772}}, we obtain \((\forall a_{1}\, :\, \alpha,\; a_{1} \in N^{-}_{r,\allowbreak{} s}(b) \implies a = a_{1}) = (\forall a_{1}\, :\, \alpha,\; a_{1} \in s \land r(a_{1},\allowbreak{} b) \implies a = a_{1})\).
By \paperlink{premise:696d706c6965735f636f6e6772}{Prerequisite~P\ref{premise:696d706c6965735f636f6e6772}} with \(p_{1} = a \in N^{-}_{r, s}(b)\) and \(p_{2} = a \in s \land r\left(a, b\right)\) and \(q_{1} = \forall a_{1}\, :\, \alpha,\; a_{1} \in N^{-}_{r, s}(b) \implies a = a_{1}\) and \(q_{2} = \forall a_{1}\, :\, \alpha,\; a_{1} \in s \land r\left(a_{1}, b\right) \implies a = a_{1}\), we obtain \((a \in N^{-}_{r,\allowbreak{} s}(b) \implies (\forall a_{2}\, :\, \alpha,\; a_{2} \in N^{-}_{r,\allowbreak{} s}(b) \implies a = a_{2})) = (a \in s \land r(a,\allowbreak{} b) \implies (\forall a_{2}\, :\, \alpha,\; a_{2} \in s \land r(a_{2},\allowbreak{} b) \implies a = a_{2}))\).
\end{proofstep}
Use the hypothesis \(b \in t\).
Applying the hypothesis \(\forall b\, :\, \beta,\; b \in t \implies \operatorname{Subsingleton}(\{a_{1}\in \alpha \mid a_{1} \in s \land r(a_{1},\allowbreak{} b)\})\), we obtain \(\forall a\, :\, \alpha,\; a \in s \land r(a,\allowbreak{} b) \implies (\forall b_{1}\, :\, \alpha,\; b_{1} \in s \land r(b_{1},\allowbreak{} b) \implies a = b_{1})\).
Rewriting with \paperlink{premise:666f72616c6c5f636f6e6772}{Prerequisite~P\ref{premise:666f72616c6c5f636f6e6772}} and the identity established above in the preceding result gives \(\forall a\, :\, \alpha,\; a \in N^{-}_{r,\allowbreak{} s}(b) \implies (\forall b_{1}\, :\, \alpha,\; b_{1} \in N^{-}_{r,\allowbreak{} s}(b) \implies a = b_{1})\).
By the reverse implication, we obtain \(\left|N^{-}_{r, s}(b)\right| \le 1\).
\end{proofstep}
By \paperlink{theorem:16}{Corollary~\ref{theorem:16}}, we obtain \(\left|s\right| \cdot 1 \le \left|t\right| \cdot 1\).
Rewriting with \paperlink{premise:6d756c5f6f6e65}{Prerequisite~P\ref{premise:6d756c5f6f6e65}} in the preceding result gives \(\left|s\right| \cdot 1 \le \left|t\right|\).
Rewriting with \paperlink{premise:6d756c5f6f6e65}{Prerequisite~P\ref{premise:6d756c5f6f6e65}} in the preceding result gives \(\left|s\right| \le \left|t\right|\).
\end{proof}

\begin{corollary}[{Bounding nonempty columns by rows of degree at most one}]\label{theorem:20}
Let \(\alpha\) and \(\beta\) be types. Let \(r\) be a relation from \(\alpha\) to \(\beta\). Let \(s\) be a finite subset of \(\alpha\). Let \(t\) be a finite subset of \(\beta\). Assume \(\forall b\, :\, \beta,\; b \in t \implies (\exists a_{1}\, :\, \alpha,\; a_{1} \in s \land r(a_{1},\allowbreak{} b))\). Assume \(\forall a\, :\, \alpha,\; a \in s \implies \operatorname{Subsingleton}(\{b\in \beta \mid b \in t \land r(a,\allowbreak{} b)\})\). Then \(\left|t\right| \le \left|s\right|\).
\end{corollary}

\begin{proof}
Apply \paperlink{theorem:19}{Lemma~\ref{theorem:19}} with \(\operatorname{swap}\left(r\right)\) in place of its parameter \(r\).
\end{proof}

\section{Finite types}
\begin{theorem}[{A cardinality bound from a left-total relation}]\label{theorem:21}
Let \(\alpha\) and \(\beta\) be types. Fix a finite enumeration of \(\alpha\). Fix a finite enumeration of \(\beta\). Let \(r\) be a relation from \(\alpha\) to \(\beta\). Assume \(\forall a\, :\, \alpha,\; \exists b\, :\, \beta,\; r(a, b)\). Assume \(\forall b\, :\, \beta,\; \forall a, a^{'}\, :\, \alpha,\; r(a, b) \land r(a^{'}, b) \implies a = a^{'}\). Then \(\left|\alpha\right| \le \left|\beta\right|\).
\end{theorem}

\begin{proof}
The substitution uses the following argument.
\begin{proofstep}
Fix \(a\) in \(\alpha\).
By \paperlink{premise:46696e7365742e6d656d5f756e6976}{Prerequisite~P\ref{premise:46696e7365742e6d656d5f756e6976}} with \(x = a\), we obtain \(a \in \alpha\).
By \paperlink{premise:65715f74727565}{Prerequisite~P\ref{premise:65715f74727565}}, we obtain \(\left(a \in \alpha\right) = \mathrm{True}\).
For the required implication, argue as follows.
\begin{proofstep}
Fix \(b\) in \(\beta\).
By \paperlink{premise:46696e7365742e6d656d5f756e6976}{Prerequisite~P\ref{premise:46696e7365742e6d656d5f756e6976}} with \(x = b\), we obtain \(b \in \beta\).
By \paperlink{premise:65715f74727565}{Prerequisite~P\ref{premise:65715f74727565}}, we obtain \(\left(b \in \beta\right) = \mathrm{True}\).
Applying the same function to equal arguments, we obtain \(\left(b \in \beta \land r\left(a, b\right)\right) = \left(\mathrm{True} \land r\left(a, b\right)\right)\).
By \paperlink{premise:747275655f616e64}{Prerequisite~P\ref{premise:747275655f616e64}} with \(p = r\left(a, b\right)\), we obtain \(\left(\mathrm{True} \land r\left(a, b\right)\right) = r\left(a, b\right)\).
Chaining the equalities, we obtain \(\left(b \in \beta \land r\left(a, b\right)\right) = r\left(a, b\right)\).
\end{proofstep}
By function extensionality, we obtain \(\left(x \mapsto x \in \beta \land r\left(a, x\right)\right) = r\left(a\right)\).
Applying the same function to equal arguments, we obtain \((\exists x\, :\, \beta,\; x \in \beta \land r(a,\allowbreak{} x)) = (\exists x\, :\, \beta,\; r(a)(x))\).
By \paperlink{premise:696d706c6965735f636f6e6772}{Prerequisite~P\ref{premise:696d706c6965735f636f6e6772}} with \(p_{1} = a \in \alpha\) and \(p_{2} = \mathrm{True}\) and \(q_{1} = \exists x\, :\, \beta,\; x \in \beta \land r\left(a, x\right)\) and \(q_{2} = \exists x\, :\, \beta,\; r\left(a\right)(x)\), we obtain \((a \in \alpha \implies (\exists x\, :\, \beta,\; x \in \beta \land r(a,\allowbreak{} x))) = (\mathrm{True} \implies (\exists x\, :\, \beta,\; r(a)(x)))\).
By \paperlink{premise:696e73744e6f6e656d7074794f66496e68616269746564}{Prerequisite~P\ref{premise:696e73744e6f6e656d7074794f66496e68616269746564}} with \(\mathrm{True}\) as its type parameter \(\alpha\), we obtain \(\mathrm{True}\text{ is nonempty}\).
By \paperlink{premise:666f72616c6c5f636f6e7374}{Prerequisite~P\ref{premise:666f72616c6c5f636f6e7374}} with \(\mathrm{True}\) as its type parameter \(\alpha\), we obtain \((\forall a_{1}\, :\, \mathrm{True},\; \exists b\, :\, \beta,\; r(a,\allowbreak{} b)) \iff (\exists b\, :\, \beta,\; r(a,\allowbreak{} b))\).
By propositional extensionality, we obtain \((\mathrm{True} \implies (\exists b\, :\, \beta,\; r(a,\allowbreak{} b))) = (\exists b\, :\, \beta,\; r(a,\allowbreak{} b))\).
Chaining the equalities, we obtain \((a \in \alpha \implies (\exists x\, :\, \beta,\; x \in \beta \land r(a,\allowbreak{} x))) = (\exists b\, :\, \beta,\; r(a,\allowbreak{} b))\).
\end{proofstep}
Use the hypothesis \(\forall a\, :\, \alpha,\; \exists b\, :\, \beta,\; r(a, b)\).
Rewriting with \paperlink{premise:666f72616c6c5f636f6e6772}{Prerequisite~P\ref{premise:666f72616c6c5f636f6e6772}} and the identity established above in the preceding result gives \(\forall a\, :\, \alpha,\; a \in \alpha \implies (\exists b\, :\, \beta,\; b \in \beta \land r(a,\allowbreak{} b))\).
For the required implication, argue as follows.
\begin{proofstep}
Fix \(x\) in \(\beta\).
Assume \(x \in \beta\).
Fix \(x_{2}\) in \(\alpha\).
Assume \(x_{2} \in \{a\in \alpha \mid a \in \alpha \land r\left(a, x\right)\}\).
Fix \(x_{3}\) in \(\alpha\).
Assume \(x_{3} \in \{a\in \alpha \mid a \in \alpha \land r\left(a, x\right)\}\).
Use the hypothesis \(x_{2} \in \{a\in \alpha \mid a \in \alpha \land r\left(a, x\right)\}\).
Taking the second conjunct, we obtain \(r\left(x_{2}, x\right)\).
Use the hypothesis \(x_{3} \in \{a\in \alpha \mid a \in \alpha \land r\left(a, x\right)\}\).
Taking the second conjunct, we obtain \(r\left(x_{3}, x\right)\).
Applying the hypothesis \(\forall b\, :\, \beta,\; \forall a, a^{'}\, :\, \alpha,\; r(a, b) \land r(a^{'}, b) \implies a = a^{'}\), we obtain \(x_{2} = x_{3}\).
\end{proofstep}
Apply \paperlink{theorem:19}{Lemma~\ref{theorem:19}} with \(s = \alpha\) and \(t = \beta\).
\end{proof}

\begin{theorem}[{A cardinality bound from a right-total relation}]\label{theorem:22}
Let \(\alpha\) and \(\beta\) be types. Fix a finite enumeration of \(\alpha\). Fix a finite enumeration of \(\beta\). Let \(r\) be a relation from \(\alpha\) to \(\beta\). Assume \(\forall b\, :\, \beta,\; \exists a\, :\, \alpha,\; r(a, b)\). Assume \(\forall a\, :\, \alpha,\; \forall b, b^{'}\, :\, \beta,\; r(a, b) \land r(a, b^{'}) \implies b = b^{'}\). Then \(\left|\beta\right| \le \left|\alpha\right|\).
\end{theorem}

\begin{proof}
The substitution uses the following argument.
\begin{proofstep}
Fix \(b\) in \(\beta\).
By \paperlink{premise:46696e7365742e6d656d5f756e6976}{Prerequisite~P\ref{premise:46696e7365742e6d656d5f756e6976}} with \(x = b\), we obtain \(b \in \beta\).
By \paperlink{premise:65715f74727565}{Prerequisite~P\ref{premise:65715f74727565}}, we obtain \(\left(b \in \beta\right) = \mathrm{True}\).
For the required implication, argue as follows.
\begin{proofstep}
Fix \(a\) in \(\alpha\).
By \paperlink{premise:46696e7365742e6d656d5f756e6976}{Prerequisite~P\ref{premise:46696e7365742e6d656d5f756e6976}} with \(x = a\), we obtain \(a \in \alpha\).
By \paperlink{premise:65715f74727565}{Prerequisite~P\ref{premise:65715f74727565}}, we obtain \(\left(a \in \alpha\right) = \mathrm{True}\).
Applying the same function to equal arguments, we obtain \(\left(a \in \alpha \land r\left(a, b\right)\right) = \left(\mathrm{True} \land r\left(a, b\right)\right)\).
By \paperlink{premise:747275655f616e64}{Prerequisite~P\ref{premise:747275655f616e64}} with \(p = r\left(a, b\right)\), we obtain \(\left(\mathrm{True} \land r\left(a, b\right)\right) = r\left(a, b\right)\).
Chaining the equalities, we obtain \(\left(a \in \alpha \land r\left(a, b\right)\right) = r\left(a, b\right)\).
\end{proofstep}
By function extensionality, we obtain \(\left(x \mapsto x \in \alpha \land r\left(x, b\right)\right) = \left(x \mapsto r\left(x, b\right)\right)\).
Applying the same function to equal arguments, we obtain \((\exists x\, :\, \alpha,\; x \in \alpha \land r(x,\allowbreak{} b)) = (\exists x\, :\, \alpha,\; r(x,\allowbreak{} b))\).
By \paperlink{premise:696d706c6965735f636f6e6772}{Prerequisite~P\ref{premise:696d706c6965735f636f6e6772}} with \(p_{1} = b \in \beta\) and \(p_{2} = \mathrm{True}\) and \(q_{1} = \exists x\, :\, \alpha,\; x \in \alpha \land r\left(x, b\right)\) and \(q_{2} = \exists x\, :\, \alpha,\; r\left(x, b\right)\), we obtain \((b \in \beta \implies (\exists x\, :\, \alpha,\; x \in \alpha \land r(x,\allowbreak{} b))) = (\mathrm{True} \implies (\exists x\, :\, \alpha,\; r(x,\allowbreak{} b)))\).
By \paperlink{premise:696e73744e6f6e656d7074794f66496e68616269746564}{Prerequisite~P\ref{premise:696e73744e6f6e656d7074794f66496e68616269746564}} with \(\mathrm{True}\) as its type parameter \(\alpha\), we obtain \(\mathrm{True}\text{ is nonempty}\).
By \paperlink{premise:666f72616c6c5f636f6e7374}{Prerequisite~P\ref{premise:666f72616c6c5f636f6e7374}} with \(\mathrm{True}\) as its type parameter \(\alpha\), we obtain \((\forall a\, :\, \mathrm{True},\; \exists a_{1}\, :\, \alpha,\; r(a_{1},\allowbreak{} b)) \iff (\exists a\, :\, \alpha,\; r(a,\allowbreak{} b))\).
By propositional extensionality, we obtain \((\mathrm{True} \implies (\exists a_{1}\, :\, \alpha,\; r(a_{1},\allowbreak{} b))) = (\exists a\, :\, \alpha,\; r(a,\allowbreak{} b))\).
Chaining the equalities, we obtain \((b \in \beta \implies (\exists x\, :\, \alpha,\; x \in \alpha \land r(x,\allowbreak{} b))) = (\exists a\, :\, \alpha,\; r(a,\allowbreak{} b))\).
\end{proofstep}
Use the hypothesis \(\forall b\, :\, \beta,\; \exists a\, :\, \alpha,\; r(a, b)\).
Rewriting with \paperlink{premise:666f72616c6c5f636f6e6772}{Prerequisite~P\ref{premise:666f72616c6c5f636f6e6772}} and the identity established above in the preceding result gives \(\forall b\, :\, \beta,\; b \in \beta \implies (\exists a_{1}\, :\, \alpha,\; a_{1} \in \alpha \land r(a_{1},\allowbreak{} b))\).
For the required implication, argue as follows.
\begin{proofstep}
Fix \(x\) in \(\alpha\).
Assume \(x \in \alpha\).
Fix \(x_{2}\) in \(\beta\).
Assume \(x_{2} \in \{b\in \beta \mid b \in \beta \land r\left(x, b\right)\}\).
Fix \(x_{3}\) in \(\beta\).
Assume \(x_{3} \in \{b\in \beta \mid b \in \beta \land r\left(x, b\right)\}\).
Use the hypothesis \(x_{2} \in \{b\in \beta \mid b \in \beta \land r\left(x, b\right)\}\).
Taking the second conjunct, we obtain \(r\left(x, x_{2}\right)\).
Use the hypothesis \(x_{3} \in \{b\in \beta \mid b \in \beta \land r\left(x, b\right)\}\).
Taking the second conjunct, we obtain \(r\left(x, x_{3}\right)\).
Applying the hypothesis \(\forall a\, :\, \alpha,\; \forall b, b^{'}\, :\, \beta,\; r(a, b) \land r(a, b^{'}) \implies b = b^{'}\), we obtain \(x_{2} = x_{3}\).
\end{proofstep}
Apply \paperlink{theorem:20}{Corollary~\ref{theorem:20}} with \(s = \alpha\) and \(t = \beta\).
\end{proof}

\end{document}
